\newpage
\subsection{Организация виртуальной вычислительной среды}
\subsection*{Выбор и развёртывание среды виртуализации}
В качестве среды виртуализации в данной работе использовался
программный продукт Oracle VM VirtualBox, выбор которого обусловлен
его свободным распространением, поддержкой широкого спектра
гостевых операционных систем и наличием официальной сборки для
архитектуры ARM64, что является необходимым условием для
развёртывания виртуальной среды на хост-системе с процессором Apple
Silicon (MacBook Air, чип Apple M2, 8 ГБ ОЗУ).

В соответствии с вариантом задания в качестве гостевой операционной
системы был выбран дистрибутив Fedora Server в редакции для
архитектуры aarch64. Установка и функционирование системы
осуществлялись в режиме интерфейса командной строки, без
графической оболочки.

\subsection*{Архитектура виртуального кластера}
Для имитации вычислительных ресурсов суперкомпьютера было развёрнуто
пять виртуальных машин: один управляющий узел (\textit{master}), на котором размещаются диспетчер SLURM и средства разработки,
и четыре вычислительных узла (\textit{node1-node4}), предоставляющие
ресурсы для выполнения заданий. Конфигурация виртуальных машин,
выбранная с учётом ограничений хост-системы по объёму оперативной
памяти, приведена в таблице~\ref{tab:nodes_config}.

\begin{table}[H]
\caption{Конфигурация узлов виртуального кластера}
\label{tab:nodes_config}
\centering
\begin{tabular}{|l|c|c|c|}
\hline
Узел & Роль & ЦП (ядра) & ОЗУ (МБ) \\
\hline
master & управляющий узел & 2 & 2048 \\
node1 & вычислительный узел & 1 & 1024 \\
node2 & вычислительный узел & 1 & 1024 \\
node3 & вычислительный узел & 1 & 1024 \\
node4 & вычислительный узел & 1 & 1024 \\
\hline
\end{tabular}
\end{table}

На каждом узле создана единая непривилегированная учётная запись
пользователя(\textit{vlad}), используемая для запуска заданий и работы с диспетчером
SLURM; учётная запись \texttt{root} применялась исключительно для
установки и настройки системного программного обеспечения.

\subsection*{Настройка виртуальной сети и разрешения имён}
Для обеспечения сетевого взаимодействия внутри кластера на каждой
виртуальной машине был настроен сетевой адаптер в режиме внутренней
сети (\texttt{Internal Network}), формирующий изолированную локальную
подсеть \texttt{192.168.1.0/24}. Для обеспечения доступа с
хост-машины к управляющему узлу дополнительно задействован интерфейс
\texttt{NAT} с пробросом порта 2222 хоста на порт 22 гостевой ОС.
Каждому узлу присвоен статический IP-адрес согласно
таблице~\ref{tab:ip}.

\begin{table}[H]
\caption{Адресация узлов виртуального кластера}
\label{tab:ip}
\centering
\begin{tabular}{|l|c|}
\hline
Имя узла & IP-адрес \\
\hline
master & 192.168.1.10 \\
node1  & 192.168.1.11 \\
node2  & 192.168.1.12 \\
node3  & 192.168.1.13 \\
node4  & 192.168.1.14 \\
\hline
\end{tabular}
\end{table}

Разрешение имён узлов реализовано посредством редактирования файла
\texttt{/etc/hosts} на каждой виртуальной машине. Данный подход
позволяет обращаться к узлам по символьным именам
(\texttt{master}, \texttt{node1}--\texttt{node4}) без развёртывания
внешнего DNS-сервера и обеспечивает корректную работу демонов
планировщика SLURM. Итоговое содержимое файла приведено в
листинге~1.

\begin{lstlisting}[caption={Содержимое файла /etc/hosts}]
127.0.0.1    localhost
192.168.1.10 master
192.168.1.11 node1
192.168.1.12 node2
192.168.1.13 node3
192.168.1.14 node4
\end{lstlisting}

\subsection*{Настройка SSH с авторизацией по ключам}
На управляющем узле для непривилегированной учётной записи
сгенерирована пара криптографических ключей --- приватный и
публичный --- с использованием утилиты \texttt{ssh-keygen}
(листинг~2). Публичный ключ распространён на все вычислительные
узлы посредством утилиты \texttt{ssh-copy-id}, что обеспечивает
выполнение удалённых подключений и команд без запроса пароля.

\begin{lstlisting}[caption={Генерация и распространение ключей SSH}]
ssh-keygen -t ed25519 -C "cluster-key"
ssh-copy-id vlad@node1
ssh-copy-id vlad@node2
ssh-copy-id vlad@node3
ssh-copy-id vlad@node4
\end{lstlisting}

В конфигурационном файле демона SSH (\texttt{/etc/ssh/sshd\_config})
на всех узлах отключена авторизация по паролю посредством установки
параметра \texttt{PasswordAuthentication~no}, что исключает
возможность входа без криптографического ключа.

\subsection*{Установка и настройка munge}
Для аутентификации сообщений, которыми обмениваются демоны SLURM на
разных узлах, на управляющем узле был сгенерирован общий
секретный ключ \texttt{munge.key}, который был скопирован на все
вычислительные узлы с сохранением единых владельца и прав доступа
(пользователь \texttt{munge}). Служба \texttt{munge} была запущена и
добавлена в автозагрузку на всех пяти узлах кластера.

\subsection*{Установка и настройка SLURM}
На управляющем узле был установлен и настроен демон диспетчера
\texttt{slurmctld}, на вычислительных узлах — демон
\texttt{slurmd}. Основные параметры кластера определены в общем для
всех узлов конфигурационном файле \texttt{/etc/slurm/slurm.conf}:
имя кластера, адрес управляющего узла, список вычислительных узлов с
указанием количества процессоров и объёма оперативной памяти, а
также определение единственного раздела (\textit{partition}),
включающего все четыре вычислительных узла.

\begin{lstlisting}[caption={Фрагмент файла slurm.conf}]
ClusterName=mac_cluster
SlurmctldHost=master
SlurmUser=slurm
SlurmdUser=root
AuthType=auth/munge
CryptoType=crypto/munge
SelectType=select/cons_tres
SelectTypeParameters=CR_Core
SlurmdParameters=config_overrides

# Логи
SlurmctldLogFile=/var/log/slurm/slurmctld.log
SlurmdLogFile=/var/log/slurm/slurmd.log
StateSaveLocation=/var/spool/slurmctld

# Описание машин
NodeName=node[1-4] CPUs=1 State=UNKNOWN
PartitionName=debug Nodes=node[1-4] Default=YES MaxTime=INFINITE State=UP
\end{lstlisting}

После запуска демонов на всех узлах корректность конфигурации и
доступность вычислительных узлов проверялась командами
\texttt{sinfo} и \texttt{scontrol show nodes}.

\begin{lstlisting}[caption={Вывод команды \texttt{sinfo}}]
vlad@master:~$ sinfo
PARTITION AVAIL  TIMELIMIT  NODES  STATE NODELIST
debug*       up   infinite      4   idle node[1-4]
\end{lstlisting}

\subsection*{Клонирование вычислительных узлов}

После полной настройки и проверки работоспособности первого
вычислительного узла (установка ОС, сеть, SSH, munge, slurmd) данная
виртуальная машина использовалась как базовый образ для клонирования
оставшихся четырех вычислительных узлов средствами VirtualBox.